% ============================================================
%  DobleA LAB — Módulo 2
%  Diseño de instrumentos y calidad en la medición
%  Miércoles 13 de mayo · 6–9 PM
% ============================================================
\documentclass[aspectratio=169, 12pt]{beamer}

% ── Encoding & language ────────────────────────────────────
\usepackage[utf8]{inputenc}
\usepackage[T1]{fontenc}
\usepackage[spanish]{babel}
\usepackage{microtype}

% ── Fonts ──────────────────────────────────────────────────
\usepackage{lmodern}

% ── Colors (DobleA palette) ────────────────────────────────
\usepackage{xcolor}
\definecolor{daPurpleDeep}{HTML}{3B0F8C}
\definecolor{daPurpleMid}{HTML}{6B21C8}
\definecolor{daPurpleLight}{HTML}{8B5CF6}
\definecolor{daMagenta}{HTML}{E040A0}
\definecolor{daDark}{HTML}{0A0A0A}
\definecolor{daLightBg}{HTML}{F0EDE8}
\definecolor{daTextMid}{HTML}{444444}

% ── Beamer theme base ──────────────────────────────────────
\usetheme{default}
\usecolortheme{default}
\setbeamertemplate{navigation symbols}{}
\setbeamertemplate{itemize item}{\color{daMagenta}\textbullet}
\setbeamertemplate{itemize subitem}{\color{daPurpleLight}$\circ$}

% ── Background ─────────────────────────────────────────────
\setbeamercolor{background canvas}{bg=daDark}
\setbeamercolor{normal text}{fg=white}
\setbeamercolor{frametitle}{fg=white, bg=daPurpleDeep}
\setbeamercolor{title}{fg=white}
\setbeamercolor{subtitle}{fg=daPurpleLight}
\setbeamercolor{author}{fg=daLightBg}
\setbeamercolor{date}{fg=daMagenta}
\setbeamercolor{institute}{fg=daLightBg}
\setbeamercolor{block title}{bg=daPurpleMid, fg=white}
\setbeamercolor{block body}{bg=daPurpleDeep!40!black, fg=white}
\setbeamercolor{alerted text}{fg=daMagenta}

% ── Frametitle style ───────────────────────────────────────
\setbeamertemplate{frametitle}{%
  \nointerlineskip
  \begin{beamercolorbox}[wd=\paperwidth, ht=0.55cm, dp=0.25cm, leftskip=0.8cm]{frametitle}%
    \usebeamerfont{frametitle}\insertframetitle%
    \ifx\insertframesubtitle\empty\else
      \hspace{0.5em}{\color{daMagenta}\small\textit{— \insertframesubtitle}}%
    \fi
  \end{beamercolorbox}%
  \vspace{-0.15cm}%
  {\color{daMagenta}\rule{\paperwidth}{1.2pt}}%
}

% ── Footer ─────────────────────────────────────────────────
\setbeamertemplate{footline}{%
  {\color{daPurpleDeep!60!black}\rule{\paperwidth}{0.5pt}}%
  \begin{beamercolorbox}[wd=\paperwidth, ht=0.45cm, dp=0.15cm]{normal text}%
    \hspace{0.6cm}%
    {\color{daMagenta}\tiny DobleA LAB}%
    \hfill%
    {\color{daLightBg!60!black}\tiny Módulo 2 · Diseño de instrumentos y calidad en la medición}%
    \hfill%
    {\color{daPurpleLight}\tiny\insertframenumber/\inserttotalframenumber\hspace{0.5cm}}%
  \end{beamercolorbox}%
}

% ── Title page template ────────────────────────────────────
\setbeamertemplate{title page}{%
  \begin{tikzpicture}[remember picture, overlay]
    % Gradient background strip
    \fill[daPurpleDeep] (current page.north west) rectangle ([yshift=-0.45\paperheight]current page.north east);
    \fill[daPurpleMid!60!daDark] ([yshift=-0.45\paperheight]current page.north west) rectangle (current page.south east);
    % Magenta accent line
    \draw[daMagenta, line width=2pt] ([yshift=-0.45\paperheight]current page.north west) -- ([yshift=-0.45\paperheight]current page.north east);
  \end{tikzpicture}%
  \vspace{0.8cm}
  \begin{center}
    {\color{daMagenta}\footnotesize\textbf{\MakeUppercase{DobleA LAB · Taller de Investigación Aplicada}}}\\[0.3cm]
    {\color{white}\LARGE\textbf{\inserttitle}}\\[0.3cm]
    {\color{daPurpleLight}\small\insertsubtitle}\\[0.6cm]
    {\color{daLightBg}\footnotesize\insertauthor}\\[0.15cm]
    {\color{daMagenta}\footnotesize\insertdate}
  \end{center}
}

% ── TikZ for title bg ──────────────────────────────────────
\usepackage{tikz}

% ── Metadata ───────────────────────────────────────────────
\title{Diseño de instrumentos\\y calidad en la medición}
\subtitle{Módulo 2}
\author{Pablo Argote \and Daniela Urbina \and Giancarlo Visconti \textit{(invitado)}}
\date{Miércoles 13 de mayo · 6--9 PM}
\institute{DobleA LAB}

% ── Helper command: session header ─────────────────────────
\newcommand{\sesion}[2]{%
  \begin{beamercolorbox}[wd=\linewidth, ht=0.4cm, dp=0.15cm, leftskip=0.3cm, rounded=true]{block title}%
    {\small\textbf{Sesión #1} \hfill \textit{\footnotesize #2}}%
  \end{beamercolorbox}\vspace{0.15cm}%
}

% ============================================================
\begin{document}

% ── Title slide ────────────────────────────────────────────
\begin{frame}[plain]
  \titlepage
\end{frame}

% ── Overview ───────────────────────────────────────────────
\begin{frame}{Contenidos del módulo}
  \begin{columns}[T]
    \begin{column}{0.32\textwidth}
      \sesion{1}{50 min}
      {\small
      \begin{itemize}
        \item De conceptos a preguntas claras
        \item Tipos de preguntas y escalas
        \item Estructura de instrumentos
      \end{itemize}}
    \end{column}
    \begin{column}{0.32\textwidth}
      \sesion{2}{50 min}
      {\small
      \begin{itemize}
        \item Diseño de cuestionarios y pautas
        \item Errores y sesgos de medición
      \end{itemize}}
    \end{column}
    \begin{column}{0.32\textwidth}
      \sesion{3}{50 min}
      {\small
      \begin{itemize}
        \item Experimentos en encuestas
        \item Diseño experimental aplicado
        \item Ejercicio práctico
      \end{itemize}}
    \end{column}
  \end{columns}
\end{frame}

% ============================================================
%  SESIÓN 1
% ============================================================
\begin{frame}[plain]
  \begin{tikzpicture}[remember picture, overlay]
    \fill[daPurpleMid] (current page.north west) rectangle (current page.south east);
    \node[anchor=center] at (current page.center) {%
      \color{white}\Huge\textbf{Sesión 1}%
    };
    \node[anchor=center, yshift=-1.2cm] at (current page.center) {%
      \color{daMagenta}\large De conceptos a preguntas medibles%
    };
  \end{tikzpicture}
\end{frame}

\begin{frame}{De la operacionalización al instrumento}
  \begin{block}{Recapitulación del Módulo 1}
    Ya tenemos una pregunta de investigación, conceptos definidos y una estrategia metodológica. Ahora hay que \alert{construir el instrumento} que recoge los datos.
  \end{block}
  \vspace{0.4cm}
  \begin{itemize}
    \item \textbf{Concepto} → \textbf{Dimensión} → \textbf{Indicador} → \textbf{Pregunta concreta}
    \item Cada pregunta del instrumento debe poder vincularse a un objetivo del estudio.
    \item Si una pregunta no responde a ningún objetivo, sobra.
  \end{itemize}
\end{frame}

\begin{frame}{Tipos de preguntas}
  \begin{columns}[T]
    \begin{column}{0.48\textwidth}
      \textbf{\color{daMagenta}Cerradas}
      \begin{itemize}
        \small
        \item Dicotómicas (sí/no)
        \item Opción múltiple (selección única)
        \item Opción múltiple (selección múltiple)
        \item Ranking / ordenamiento
        \item Escalas (Likert, diferencial semántico)
      \end{itemize}
      \vspace{0.2cm}
      {\color{daTextMid}\tiny Ventaja: comparables, codificables, rápidas.}
    \end{column}
    \begin{column}{0.48\textwidth}
      \textbf{\color{daPurpleLight}Abiertas}
      \begin{itemize}
        \small
        \item Respuesta libre corta
        \item Respuesta libre extendida
        \item Asociación espontánea
        \item Completación de frases
      \end{itemize}
      \vspace{0.2cm}
      {\color{daTextMid}\tiny Ventaja: riqueza, profundidad, descubrimiento.}
    \end{column}
  \end{columns}
  \vspace{0.4cm}
  {\small La elección depende de si queremos \alert{medir} o \alert{explorar}.}
\end{frame}

\begin{frame}{Escalas de medición}
  \begin{itemize}
    \item \textbf{Likert:} grado de acuerdo (1 = muy en desacuerdo … 5 = muy de acuerdo).
    \begin{itemize}
      \item Decidir: ¿punto neutro o escala par (sin neutro)?
    \end{itemize}
    \item \textbf{Diferencial semántico:} pares de adjetivos opuestos (moderno–anticuado, confiable–poco confiable).
    \item \textbf{NPS (Net Promoter Score):} 0--10, clasifica en detractores · pasivos · promotores.
    \item \textbf{Escalas visuales:} termómetros, estrellas, emojis. Útiles en contextos de baja escolaridad o mobile.
  \end{itemize}
  \vspace{0.3cm}
  \begin{alertblock}{Cuidado}
    La escala define el nivel de medición (ordinal vs.\ intervalar) y condiciona qué análisis puedes hacer después.
  \end{alertblock}
\end{frame}

\begin{frame}{Estructura de un instrumento}
  \begin{enumerate}
    \item \textbf{Presentación:} quién hace el estudio, propósito, confidencialidad, duración estimada.
    \item \textbf{Filtro / screening:} preguntas para verificar que el participante califica.
    \item \textbf{Cuerpo temático:} bloques organizados por tema; de lo general a lo específico.
    \item \textbf{Datos de clasificación:} sociodemográficos al final (edad, ingreso, educación).
    \item \textbf{Cierre:} agradecimiento, espacio para comentarios, datos de contacto.
  \end{enumerate}
  \vspace{0.3cm}
  \begin{block}{Principio de flujo}
    El instrumento debe sentirse como una conversación lógica, no como un interrogatorio.
  \end{block}
\end{frame}

% ── Descanso ───────────────────────────────────────────────
\begin{frame}[plain]
  \begin{tikzpicture}[remember picture, overlay]
    \fill[daDark] (current page.north west) rectangle (current page.south east);
    \draw[daMagenta, line width=1pt] ([xshift=2cm, yshift=-2cm]current page.north west) rectangle ([xshift=-2cm, yshift=2cm]current page.south east);
    \node[anchor=center] at (current page.center) {%
      \color{daPurpleLight}\large\textit{— break 10 min —}%
    };
  \end{tikzpicture}
\end{frame}

% ============================================================
%  SESIÓN 2
% ============================================================
\begin{frame}[plain]
  \begin{tikzpicture}[remember picture, overlay]
    \fill[daPurpleMid] (current page.north west) rectangle (current page.south east);
    \node[anchor=center] at (current page.center) {%
      \color{white}\Huge\textbf{Sesión 2}%
    };
    \node[anchor=center, yshift=-1.2cm] at (current page.center) {%
      \color{daMagenta}\large Diseño de cuestionarios y sesgos%
    };
  \end{tikzpicture}
\end{frame}

\begin{frame}{Diseño de cuestionarios}
  \framesubtitle{Buenas prácticas}
  \begin{itemize}
    \item \textbf{Una idea por pregunta.} Evitar preguntas dobles:\\
          {\small\color{daTextMid} Mal: ``¿El servicio fue rápido y amable?'' → ¿qué pasa si fue rápido pero no amable?}
    \item \textbf{Lenguaje simple y neutro.} Evitar jerga técnica, dobles negaciones y palabras cargadas.
    \item \textbf{Opciones exhaustivas y excluyentes.} Que cubran todas las posibilidades sin superponerse.
    \item \textbf{Orden lógico.} Primero lo general, luego lo específico; lo menos sensible antes que lo más sensible.
    \item \textbf{Longitud razonable.} Cada pregunta adicional tiene un costo en calidad de respuesta.
  \end{itemize}
\end{frame}

\begin{frame}{Diseño de pautas cualitativas}
  \begin{columns}[T]
    \begin{column}{0.48\textwidth}
      \textbf{\color{daMagenta}Pauta de entrevista}
      \begin{itemize}
        \small
        \item Guía temática, no cuestionario rígido.
        \item Preguntas abiertas que inviten al relato.
        \item Sondeos preparados (``cuéntame más'', ``¿a qué te refieres?'').
        \item Secuencia: rapport → temas centrales → cierre reflexivo.
      \end{itemize}
    \end{column}
    \begin{column}{0.48\textwidth}
      \textbf{\color{daPurpleLight}Pauta de focus group}
      \begin{itemize}
        \small
        \item Menos preguntas, más interacción.
        \item Incluir estímulos (imágenes, conceptos, prototipos).
        \item Pregunta de apertura fácil para romper el hielo.
        \item Pregunta de cierre: ``¿hay algo que no hayamos tocado?''
      \end{itemize}
    \end{column}
  \end{columns}
  \vspace{0.3cm}
  {\small En ambos casos: \alert{pilotar} antes de salir a terreno.}
\end{frame}

\begin{frame}{Errores y sesgos de medición}
  \begin{itemize}
    \item \textbf{Deseabilidad social:} el participante responde lo que cree ``correcto'' o socialmente aceptable.
    \begin{itemize}
      \item {\small Ej.: sobrereportar reciclaje, subreportar consumo de alcohol.}
    \end{itemize}
    \item \textbf{Aquiescencia:} tendencia a responder ``sí'' o ``de acuerdo'' a todo.
    \begin{itemize}
      \item {\small Solución parcial: incluir ítems invertidos.}
    \end{itemize}
    \item \textbf{Efecto de orden:} la respuesta cambia según el lugar de la pregunta en el instrumento.
    \item \textbf{Efecto de anclaje:} la primera opción o valor presentado condiciona la respuesta.
    \item \textbf{Fatiga del encuestado:} las respuestas se degradan conforme avanza el instrumento.
  \end{itemize}
\end{frame}

\begin{frame}{¿Cómo mitigar sesgos?}
  \begin{columns}[T]
    \begin{column}{0.48\textwidth}
      \textbf{\color{daMagenta}En el diseño}
      \begin{itemize}
        \small
        \item Redacción neutra y balanceada.
        \item Aleatorizar orden de opciones.
        \item Incluir ítems de control e invertidos.
        \item Garantizar anonimato real (y comunicarlo).
        \item Usar formatos indirectos para temas sensibles.
      \end{itemize}
    \end{column}
    \begin{column}{0.48\textwidth}
      \textbf{\color{daPurpleLight}En la aplicación}
      \begin{itemize}
        \small
        \item Pilotaje con participantes reales.
        \item Entrevista cognitiva: pedir al participante que ``piense en voz alta'' mientras responde.
        \item Monitorear tiempos de respuesta.
        \item Revisar patrones sospechosos (straightlining, respuestas contradictorias).
      \end{itemize}
    \end{column}
  \end{columns}
\end{frame}

% ── Descanso ───────────────────────────────────────────────
\begin{frame}[plain]
  \begin{tikzpicture}[remember picture, overlay]
    \fill[daDark] (current page.north west) rectangle (current page.south east);
    \draw[daMagenta, line width=1pt] ([xshift=2cm, yshift=-2cm]current page.north west) rectangle ([xshift=-2cm, yshift=2cm]current page.south east);
    \node[anchor=center] at (current page.center) {%
      \color{daPurpleLight}\large\textit{— break 10 min —}%
    };
  \end{tikzpicture}
\end{frame}

% ============================================================
%  SESIÓN 3
% ============================================================
\begin{frame}[plain]
  \begin{tikzpicture}[remember picture, overlay]
    \fill[daPurpleMid] (current page.north west) rectangle (current page.south east);
    \node[anchor=center] at (current page.center) {%
      \color{white}\Huge\textbf{Sesión 3}%
    };
    \node[anchor=center, yshift=-1.2cm] at (current page.center) {%
      \color{daMagenta}\large Experimentos en encuestas%
    };
  \end{tikzpicture}
\end{frame}

\begin{frame}{¿Por qué experimentar dentro de una encuesta?}
  \begin{block}{El problema}
    A veces preguntar directamente no funciona: el participante no sabe, no quiere o no puede responder con precisión.
  \end{block}
  \vspace{0.4cm}
  \begin{itemize}
    \item Los \alert{experimentos de encuesta} (survey experiments) asignan aleatoriamente a los participantes a versiones distintas de una pregunta o estímulo.
    \item Permiten estimar el \textbf{efecto causal} de un cambio en el fraseo, en la información disponible o en el contexto.
    \item Son una herramienta para \textbf{medir mejor} lo que las preguntas directas no capturan bien.
  \end{itemize}
\end{frame}

\begin{frame}{Tipos de diseños experimentales en encuestas}
  \begin{itemize}
    \item \textbf{Experimento de viñeta (vignette):} se presenta un escenario que varía aleatoriamente en uno o más atributos.
    \begin{itemize}
      \item {\small Ej.: evaluar un candidato cuyo CV varía en género, universidad, experiencia.}
    \end{itemize}
    \item \textbf{Experimento de lista (list experiment):} técnica para medir prevalencia de conductas sensibles sin que el encuestado revele su respuesta individual.
    \item \textbf{Experimento de conjoint:} se presentan perfiles con múltiples atributos combinados; el participante elige o evalúa.
    \begin{itemize}
      \item {\small Ej.: ¿qué atributos pesan más al elegir un plan de salud?}
    \end{itemize}
    \item \textbf{Framing / priming:} variar el encuadre de una pregunta para medir cómo afecta la respuesta.
  \end{itemize}
\end{frame}

\begin{frame}{Ejemplo: experimento de viñeta}
  \begin{columns}[T]
    \begin{column}{0.48\textwidth}
      \begin{block}{\centering Grupo A}
        \small``María, 32 años, egresada de una universidad pública, con 5 años de experiencia en el sector. ¿Qué tan apta la consideras para el cargo?''
      \end{block}
    \end{column}
    \begin{column}{0.48\textwidth}
      \begin{block}{\centering Grupo B}
        \small``Pedro, 32 años, egresado de una universidad privada, con 5 años de experiencia en el sector. ¿Qué tan apto lo consideras para el cargo?''
      \end{block}
    \end{column}
  \end{columns}
  \vspace{0.4cm}
  \begin{itemize}
    \item La asignación aleatoria permite aislar el efecto de \alert{género} y \alert{tipo de universidad}.
    \item El participante no sabe que hay otra versión → reduce deseabilidad social.
    \item La diferencia entre grupos es el \textbf{efecto causal estimado}.
  \end{itemize}
\end{frame}

\begin{frame}{Claves para un buen experimento de encuesta}
  \begin{enumerate}
    \item \textbf{Aleatorización efectiva:} que los grupos sean comparables en todo excepto el tratamiento.
    \item \textbf{Tamaño de muestra suficiente:} el experimento requiere \textit{n} por grupo, no solo \textit{n} total.
    \item \textbf{Estímulo realista:} la viñeta o el escenario debe ser creíble para el participante.
    \item \textbf{Una sola variación a la vez} (o diseño factorial si se quieren interacciones).
    \item \textbf{Pre-registro:} definir hipótesis y análisis antes de recolectar datos.
  \end{enumerate}
  \vspace{0.3cm}
  \begin{alertblock}{Ventaja clave}
    Combinan la \textbf{representatividad} de la encuesta con la \textbf{lógica causal} del experimento.
  \end{alertblock}
\end{frame}

\begin{frame}{Ejercicio práctico}
  \begin{block}{Instrucción}
    En grupos de 3, elijan uno de los siguientes escenarios y diseñen un instrumento breve:
  \end{block}
  \vspace{0.3cm}
  \textbf{Escenario A:} Una empresa de delivery quiere entender por qué las evaluaciones de sus repartidores son más bajas en comunas de altos ingresos.\\[0.3cm]
  \textbf{Escenario B:} Una fundación educacional quiere medir la percepción de calidad de sus talleres entre apoderados.\\[0.3cm]
  \begin{enumerate}
    \item Redacta \textbf{5 preguntas} (al menos 2 cerradas y 1 abierta).
    \item Diseña \textbf{1 experimento de viñeta} que mida algo que una pregunta directa no capturaría bien.
    \item Identifica \textbf{2 sesgos potenciales} de tu instrumento y propón cómo mitigarlos.
  \end{enumerate}
\end{frame}

% ── Cierre ─────────────────────────────────────────────────
\begin{frame}{Cierre del módulo}
  \begin{columns}[T]
    \begin{column}{0.48\textwidth}
      \textbf{\color{daMagenta}Lo que aprendimos hoy}
      \begin{itemize}
        \small
        \item Traducir conceptos en preguntas claras.
        \item Tipos de preguntas, escalas y estructura.
        \item Diseñar cuestionarios y pautas cualitativas.
        \item Identificar y mitigar sesgos de medición.
        \item Usar experimentos para medir mejor.
      \end{itemize}
    \end{column}
    \begin{column}{0.48\textwidth}
      \textbf{\color{daPurpleLight}Próximo módulo}
      \begin{itemize}
        \small
        \item Trabajo de campo y recolección de datos.
        \item Gestión de equipos en terreno.
        \item Miércoles 20 de mayo · 6--9 PM
      \end{itemize}
    \end{column}
  \end{columns}
  \vspace{0.5cm}
  \begin{center}
    {\color{daMagenta}\rule{6cm}{0.8pt}}\\[0.3cm]
    {\small\color{daLightBg!70!black} Preguntas · Comentarios · Dudas}
  \end{center}
\end{frame}

\end{document}
